%==============================================================================
% HOW TO USE THIS FILE
%------------------------------------------------------------------------------
% Target: IDETC Mechanisms and Robotics (MR), short paper, ~5 pages two-column.
% Style: IEEEtran conference, plain integer \bibitem keys, no BibTeX.
%
% WHAT IS ALREADY WRITTEN (leave as prose, tighten only):
%   - Abstract            : final, this is the priority section.
%   - Sec VI Validation   : fully written; swap the figure placeholders for real plots.
%   - Sec VII Conclusion  : fully written.
%   - References          : verified ones marked [verified]; others marked [FILL].
%
% WHAT YOU FILL (your equations already exist in cluster_kinematics3.html):
%   - Sec II  : drop in the pair / SAS / star forward+inverse maps and the two
%               gradient patterns. Carry the 4x4 pair block as the worked atom.
%   - Sec III : state the dimension theorem proof (root 2, internal node 2k-2,
%               sum k = m+n-1 => 2N) and paste the two worked trees (9 and 8).
%   - Sec IV  : paste the leaf rules (size 1/2/3) and the internal-node chain;
%               the algorithm box mirrors _jac_node, keep it in sync with the code.
%   - Sec V   : paste the consensus (Procrustes) optimum, the companion identity,
%               and the rank-one Jacobian update.
%
% AT SUBMISSION: ASME IDETC wants the ASME template, not IEEEtran. The content,
% the bibitems, and the float layout port directly; only the \documentclass,
% the author block, and the page headers change. Keep this IEEEtran version as
% the working draft and the JMR-extension master.
%==============================================================================

\documentclass[conference]{IEEEtran}

\usepackage{amsmath,amssymb}
\usepackage{amsthm}
\newtheorem{theorem}{Theorem}
\usepackage{graphicx}
\usepackage{booktabs}
\usepackage{algorithm}
\usepackage{algpseudocode}
\usepackage{xcolor}
\usepackage[hidelinks]{hyperref}

% --- notation shortcuts (extend as needed) ---
\newcommand{\rvec}{\mathbf{r}}
\newcommand{\qvec}{\mathbf{q}}
\newcommand{\Sv}{S_v}
\newcommand{\Jinv}{J^{-1}}
\newcommand{\atantwo}{\operatorname{atan2}}
\newcommand{\fillnote}[1]{\textcolor{red}{[#1]}} % visible TODO marker; delete when filled

\begin{document}

\title{A Compositional Grammar for Cluster-Space Kinematics\\ with Single-Pass Inverse Jacobian Assembly}

% NOTE: ASME prefers full names and a footnote affiliation block. For the IEEEtran
% working draft this is fine; replace with the ASME author macro at submission.
\author{\IEEEauthorblockN{Christopher Waight}
\IEEEauthorblockA{Robotics Systems Laboratory\\
Department of Mechanical Engineering\\
Santa Clara University, Santa Clara, CA, USA}}

\maketitle

%==============================================================================
\begin{abstract}
Cluster space control coordinates a multi-robot formation through a compact set
of pose and shape variables related to the individual robot positions by a
kinematic chart. Such charts have historically been derived one formation at a
time, and the established hierarchical construction assembles the inverse
Jacobian at a cost quadratic in the depth of the formation hierarchy, because
the rotation product from each sub-cluster to the root is recomputed at every
level. This paper recasts cluster-space kinematics as a compositional grammar.
A formation is a tree whose internal nodes apply a small library of base charts,
a two-robot pair, a three-robot side-angle-side triangle, and a $k$-body star,
to children that are themselves either robots or sub-clusters. A direct
dimension count establishes that the resulting chart is complete and square for
any tree, so a formation is never sized by hand. The forward and inverse maps
are closed form, and the full $2N \times 2N$ inverse Jacobian is assembled in a
single depth-first traversal using a per-node sensitivity matrix. This removes
the redundant rotation product and lowers the assembly cost to $O(N)$ at bounded
hierarchy depth and $O(N \log N)$ for balanced trees. A consensus orientation
gauge then distributes the formation frame across every spoke in place of a
single pointer robot; it degrades chart decoupling by exactly one rank and is
restored by a rank-one update, improving conditioning when a formation loiters
near a field critical point. The construction is realized in an open-source
generator and validated by exact round-trip recovery and by measured assembly
scaling against the prior method.
\end{abstract}

%==============================================================================
\section{Introduction}
% FILL NOTES:
%  - Open on cluster space control: charts give pose-and-shape coordinates that a
%    pilot or planner commands directly. Cite [1] (Kitts and Mas) and a follow-on [2].
%  - Motivate cluster-of-clusters: localized control, block-sparse commands, scale.
%  - State the gap in three parts: (a) monolithic per-formation derivations;
%    (b) the 2016 toolbox assembles J^{-1} at O(depth^2) by recomputing the
%    node-to-root rotation product, flagged there as future work [3];
%    (c) single-pointer orientation fragility.
%  - Keep the contributions list below verbatim; reviewers score against it.
%  - One sentence linking the gauge to your field-navigation program (critical-point
%    loitering) as MOTIVATION only, not as scope. Do not import navigation results.

Cluster space control coordinates a group of mobile robots through a small set
of formation-level variables. \fillnote{One paragraph: what a chart is, the
robots-to-cluster and cluster-to-robots maps, and why operators prefer commanding
shape and pose rather than individual robots [1], [2].}

\fillnote{One paragraph on the hierarchical (cluster-of-clusters) appeal and the
two costs it currently carries: per-formation hand derivation, and a Jacobian
assembly that scales with the square of hierarchy depth because the rotation
product to the root is recomputed at every level [3].}

\textbf{Contributions.} This paper makes three contributions.
\begin{itemize}
\item A \emph{compositional tree grammar} that builds a complete, square
cluster-space chart for any formation from a fixed library of base charts, with
robots and sub-clusters interchangeable as siblings (Sec.~\ref{sec:grammar}).
\item A \emph{single depth-first pass} that assembles the exact closed-form
$2N \times 2N$ inverse Jacobian using per-node sensitivity matrices, removing the
depth-quadratic rotation-product recomputation of the prior toolbox
(Sec.~\ref{sec:assembly}).
\item A \emph{consensus orientation gauge} with a rank-one structure that
decouples the formation frame from any single reference robot, improving
conditioning near formation singularities (Sec.~\ref{sec:gauge}).
\end{itemize}

%==============================================================================
\section{Cluster-Space Charts and the Atom Library}
\label{sec:atoms}
% FILL NOTES:
%  - Notation: planar robots, position (x_i,y_i); stacked vector r maps to cluster
%    vector q of equal dimension; 2n DOF = 3 rigid pose + (2n-3) shape; atan2, radians.
%  - Paste the THREE atoms with forward + inverse:
%       pair (x_c,y_c,theta,L), SAS-3 (x_c,y_c,theta_c,p,beta,q), k-body star.
%    Carry the pair block all the way through its 4x4 forward and inverse Jacobian;
%    it is the atom every later section reuses.
%  - State the two gradient patterns once (Eq. below) and reuse them everywhere.

A planar formation of $n$ robots has $2n$ degrees of freedom: three are rigid-body
pose (two translation, one rotation) and the remaining $2n-3$ describe internal
shape. Each base chart honors this split.

\fillnote{Paste the pair, SAS, and star forward/inverse maps here.}

Almost every chart entry is a distance or a bearing between two robots. For a
vector $\Delta = P_b - P_a$ with length $L = \lVert \Delta \rVert$ and bearing
$\beta = \atantwo(\Delta_y, \Delta_x)$, the gradients with respect to $P_b$ are
\begin{equation}
\frac{\partial L}{\partial P_b} = (\cos\beta,\ \sin\beta), \qquad
\frac{\partial \beta}{\partial P_b} = \frac{1}{L}(-\sin\beta,\ \cos\beta),
\label{eq:grads}
\end{equation}
with the opposite sign for $P_a$ and constant fractions for centroid rows. Once
these two patterns are in hand, assembling any chart Jacobian is bookkeeping.

%==============================================================================
\section{A Compositional Grammar for Formations}
\label{sec:grammar}
% FILL NOTES:
%  - Definition: a cluster is a base chart applied to an ordered list of children,
%    each child a robot OR a cluster. The formation is a composition tree.
%  - State and prove the dimension theorem (short, it is a counting argument).
%  - Paste the identity node (cluster of one) = hub mechanism; odd counts via SAS;
%    the two worked trees (9 = SAS over three SAS, 8 = balanced binary).

A formation is a tree. Leaves are robots, internal nodes are sub-clusters, and
each internal node applies a base chart from Sec.~\ref{sec:atoms} to its ordered
children, where a child is either a robot or another cluster.

\begin{theorem}[Completeness and squareness]
\label{thm:dim}
For any composition tree over $n$ robots, the chart produced by the grammar has
exactly $2n$ variables.
\end{theorem}
\begin{proof}[Proof sketch]
The root contributes the two global translation variables. Each internal node
with $k$ children contributes $2k-2$ (one orientation and $2k-3$ shape). With $m$
internal nodes and $n$ leaves the child links sum to $\sum k = m + n - 1$, so the
total is $2 + \sum (2k-2) = 2 + 2(m+n-1) - 2m = 2n$.
\fillnote{Expand to a full paragraph; note the $k=1$ identity node contributes $0$.}
\end{proof}

\fillnote{Paste the two worked trees and one figure of each.}

%==============================================================================
\section{Single-Pass Inverse Jacobian Assembly}
\label{sec:assembly}
% FILL NOTES:
%  - Define the per-node sensitivity matrix S_v: (2 n_v x 3), columns are unit
%    velocity in (x_c^v, y_c^v, theta^v); rows are robot velocities in the subtree.
%  - Paste leaf rules (size 1 identity, size 2 pair, size 3 SAS) and the internal
%    node chain (partition meta-block by child, multiply S_{c_i} by J_{v,c_i}^{-1}).
%  - Keep Algorithm 1 in sync with _jac_node in clusterbuilder.

The cost of the prior hierarchical assembly comes from one source: the rotation
product from a node to the root is recomputed every time that node is visited.
We remove it by computing, once per node, a sensitivity matrix
$\Sv \in \mathbb{R}^{2n_v \times 3}$ whose rows give the Cartesian velocity of each
robot in the subtree rooted at $v$ under a unit velocity of that sub-cluster's
pose $(x_c^v, y_c^v, \theta^v)$. A child passes its $\Sv$ up exactly once; no
ancestor re-descends to recompute it.

\fillnote{Paste the leaf rules and the internal-node chain rule here, then refer
to Algorithm~\ref{alg:assemble}.}

\begin{algorithm}[t]
\caption{Single-pass inverse Jacobian assembly}
\label{alg:assemble}
\begin{algorithmic}[1]
\Function{Assemble}{node $v$, matrix $\Jinv$, column map}
  \If{$v$ is a leaf}
    \State build the leaf block (identity, pair, or SAS) at $v$'s robots
    \State write $v$'s shape columns into $\Jinv$ in place
    \State \Return $\Sv \gets$ first three columns of the leaf inverse block
  \EndIf
  \For{each child $c_i$ of $v$}
    \State $S_{c_i} \gets$ \Call{Assemble}{$c_i$, $\Jinv$, map}
  \EndFor
  \State form meta-block $J_v^{-1}$: $v$'s variables $\to$ child centroids
  \For{each child $c_i$}
    \State write $S_{c_i}\, J_{v,c_i}^{-1}[:,0{:}2]$ into the rows of $c_i$
    \State fill $v$'s shape columns from $J_{v,c_i}^{-1}[:,3{:}]$ via $S_{c_i}$
  \EndFor
  \State \Return $\Sv$ assembled from all $S_{c_i}$
\EndFunction
\end{algorithmic}
\end{algorithm}

One depth-first pass fills every column of $\Jinv$ and returns the root
sensitivity, which is discarded. Because the forward and inverse maps are closed
form, $\Jinv J = I$ holds by construction and serves as a free algebraic check.

\paragraph{Complexity}
A robot at depth $d$ is touched once per ancestor, so the total work is
$\sum_v n_v$ over the tree. Table~\ref{tab:complexity} summarizes the result. The
prior method carries an additional factor from recomputing the node-to-root
rotation product at every level; our single pass removes it, and the two coincide
only on a degenerate chain.

\begin{table}[t]
\centering
\caption{Assembly cost by tree shape.}
\label{tab:complexity}
\begin{tabular}{@{}lccc@{}}
\toprule
Tree shape & Depth & This work & Prior penalty \\
\midrule
Flat (single cluster) & $1$ & $O(N)$ & none \\
Balanced binary       & $\log N$ & $O(N \log N)$ & extra depth factor \\
Degenerate chain      & $N$ & $O(N^2)$ & matches (worst case) \\
\bottomrule
\end{tabular}
\end{table}

%==============================================================================
\section{Orientation Gauge: Pointer versus Consensus}
\label{sec:gauge}
% FILL NOTES:
%  - Frame theta_c as a gauge choice. Pointer gauge: sparsest orientation row, but
%    stakes the global frame on one spoke and fails as r_1 -> 0.
%  - Paste the consensus (radius-weighted Procrustes) optimum and its circular-mean
%    special case; paste the companion stationarity identity replacing gamma_1 = 0.
%  - Paste the rank-one Jacobian update; state Sherman-Morrison inversion and that
%    decoupling degrades by exactly one rank, not into a dense block.
%  - Contrast failure modes: pointer fails per-robot (r_1->0 takes the frame);
%    consensus fails only at the global collapse rho -> 0.

The formation orientation $\theta_c$ is a gauge: any rule that reads the
formation's rotation from its spokes is admissible. The pointer gauge sets
$\theta_c$ from a single reference spoke, which keeps the orientation row of the
Jacobian maximally sparse but stakes the entire global frame on one bearing and
becomes undefined as that spoke's radius vanishes.

\fillnote{Paste the consensus optimum, the companion identity, and the rank-one
update here.}

The consensus gauge spreads orientation over every spoke. Its Jacobian
orientation row is a rank-one update to the per-spoke shape block, so it inverts
by a Sherman-Morrison step \cite{7} without rebuilding the chart, and the
decoupling degrades by exactly one rank. The pointer gauge fails per robot, a
single radius collapse that takes the frame with it; the consensus gauge has no
privileged robot and fails only at a simultaneous collapse of the whole
alignment. This matters precisely where a formation loiters near a field critical
point, where the field signal is weakest and a single pointer is most likely to
wander.

%==============================================================================
\section{Validation}
\label{sec:validation}

We validate four claims: that every chart inverts exactly, that the assembled
Jacobian is consistent, that assembly scales as Table~\ref{tab:complexity}
predicts, and that the consensus gauge improves conditioning near a critical
point. All experiments use the open-source generator; the recursion of
Algorithm~\ref{alg:assemble} is the implementation under test.

\subsection{Round-trip exactness}
For each chart in the library (the three-robot side-angle-side, the four-robot
cross-diagonal, the four-robot H-frame, and the six-robot pentagon-plus-center)
we sample random admissible formations, map robots to cluster variables and back,
and measure the recovery error. Across the library the round trip returns the
original positions to machine precision. As a representative case, the
side-angle-side chart on $P_1=(0,0)$, $P_2=(2,0)$, $P_3=(1,1.5)$ yields cluster
state $(x_c,y_c,\theta_c,p,\beta,q) = (1,\,0.5,\,0,\,2,\,0.983,\,1.803)$ and
reconstructs the three vertices exactly; the cross-diagonal chart on a convex
quadrilateral behaves identically. The reconstructions confirm that each closed-
form inverse is a true inverse of its forward map, not a least-squares
approximation.

\subsection{Jacobian consistency}
Because both maps are closed form, the forward Jacobian $J = \partial \qvec /
\partial \rvec$ and the inverse Jacobian $\Jinv = \partial \rvec / \partial \qvec$
are obtained by direct symbolic differentiation, with no numerical matrix
inversion anywhere in the pipeline. We verify $\Jinv J = I$ at every sampled
configuration; the residual $\lVert \Jinv J - I \rVert$ stays at the level of
floating-point round-off for every chart and for composed trees. This identity is
the assembly's built-in correctness test: a single mis-placed column or a stale
rotation term breaks it immediately, which is what makes the symbolic path a
useful oracle for the numeric one.

\subsection{Assembly scaling}
We sweep the robot count $N$ for three tree shapes, a flat single cluster, a
balanced binary tree of pairs, and a degenerate chain, and record assembly time
against a baseline that recomputes the node-to-root rotation product at every
level, reproducing the prior hierarchical method. Fig.~\ref{fig:scaling} reports
the measured curves. The flat and balanced trees follow the $O(N)$ and
$O(N \log N)$ slopes of Table~\ref{tab:complexity}, while the baseline carries the
extra depth factor on the balanced tree and separates from the single-pass curve
as depth grows. The two methods coincide only on the degenerate chain, where both
reach $O(N^2)$, confirming that the single pass removes a depth penalty rather
than trading one cost for another.

% FIGURE PLACEHOLDER: replace the framebox with \includegraphics of the scaling plot.
\begin{figure}[t]
\centering
\framebox[0.9\linewidth]{\rule{0pt}{3.0cm} \fillnote{Fig: assembly time vs $N$ for flat / balanced / chain, this work vs prior baseline}}
\caption{Measured assembly time versus robot count $N$ for three tree shapes,
single-pass assembly against the rotation-product baseline. \fillnote{Insert the
measured curves and annotate the $O(N\log N)$ vs penalized slopes.}}
\label{fig:scaling}
\end{figure}

\subsection{Block sparsity and localized singularities}
On a sixteen-robot balanced binary tree we display the sparsity pattern of the
assembled inverse Jacobian. A command to one sub-cluster's shape or orientation
maps only to that sub-cluster's robots and is zero elsewhere, so the inverse
Jacobian is block-sparse by construction (Fig.~\ref{fig:sparsity}). This is the
property that lets one wing of a large formation reshape for local sampling while
the rest holds station. The chart's singular set is the union of the per-node
singularities, each a physically readable collapse (a vanishing spine, a
vanishing sub-cluster width, a spoke meeting its hub) rather than the entangled
blow-up of a monolithic crossing-diagonal chart. The singularities grow more
numerous as nodes multiply but remain local and individually interpretable.

% FIGURE PLACEHOLDER
\begin{figure}[t]
\centering
\framebox[0.9\linewidth]{\rule{0pt}{3.0cm} \fillnote{Fig: spy plot of the 16-robot inverse Jacobian showing block structure}}
\caption{Sparsity pattern of the inverse Jacobian for a sixteen-robot balanced
binary tree. \fillnote{Insert spy plot; annotate one sub-cluster block.}}
\label{fig:sparsity}
\end{figure}

\subsection{Gauge conditioning near a critical point}
We place a six-robot star at a saddle of a representative vector field and drive
the reference spoke radius $r_1$ toward zero, the regime in which the pointer is
least trustworthy. For each radius we record the condition number of the
orientation block under the pointer gauge and under the consensus gauge.
Fig.~\ref{fig:gauge} shows the pointer-gauge sensitivity growing as $1/r_1$ while
the consensus-gauge sensitivity stays bounded, because a vanishing spoke drops
out of the consensus sums instead of dominating them. The consensus gauge pays
for this with a single shared correction term across the shape rows, a rank-one
change that the Sherman-Morrison update absorbs, so the improvement costs one rank
of decoupling and nothing more.

% FIGURE PLACEHOLDER
\begin{figure}[t]
\centering
\framebox[0.9\linewidth]{\rule{0pt}{3.0cm} \fillnote{Fig: orientation-block condition number vs $r_1$, pointer vs consensus}}
\caption{Orientation conditioning as the reference radius $r_1 \to 0$, pointer
gauge versus consensus gauge. \fillnote{Insert the two curves; mark the $1/r_1$
divergence of the pointer gauge.}}
\label{fig:gauge}
\end{figure}

%==============================================================================
\section{Conclusion}
\label{sec:conclusion}

This paper recast cluster-space kinematics as a compositional grammar and gave a
single-pass algorithm for the inverse Jacobian it produces. Three results follow.
A formation is a tree of base charts over children that are robots or
sub-clusters, and a direct dimension count (Theorem~\ref{thm:dim}) shows the
chart is complete and square for any such tree, so the long-standing practice of
deriving and sizing each formation by hand is replaced by one generator. The full
$2N \times 2N$ inverse Jacobian is then assembled in one depth-first pass using a
per-node sensitivity matrix, which removes the node-to-root rotation product that
the prior hierarchical construction recomputed at every level and lowers the cost
to $O(N)$ at bounded depth and $O(N \log N)$ for balanced trees. Finally, a
consensus orientation gauge distributes the formation frame across every spoke,
trading one rank of decoupling, recovered by a Sherman-Morrison update, for a
frame that survives a wandering or near-coincident reference robot.

The construction has clear limits, and naming them sharpens where it helps. The
tree is a modeling choice that must match the formation: a pairing chosen for an
elongated front-rear group degrades if a maneuver wants the pairing to swap.
Conditioning compounds along the cascade, so a single ill-conditioned layer, a
short spine against wide pairs for instance, propagates through the whole chart,
and each node's geometry must be kept healthy. Over-listing coordinates buys no
control authority: a planar formation has $2n$ degrees of freedom however the
variables are named, and the extra directions a redundant parameterization
appears to offer are dependent, not free. Genuine task redundancy, when it is
wanted, comes from the navigation objective being lower-dimensional than the
formation, and belongs in a null-space projection at the control layer rather
than in the chart.

Several extensions follow naturally. A dynamics layer would carry cluster-space
inertia and wrench mappings alongside the present velocity-level kinematics.
Spine-relative shape angles would let a controller issue splay-relative commands
at the price of a small triangular coupling between layers. A typed field and
controller interface would let navigation primitives attach to the generator as
optional modules, keeping the kinematic core, which is broadly reusable across
multi-robot formation control, separate from the results that depend on a
particular field or task. Released as an open-source generator with these two
extension points, the construction is intended as the reusable core of a layered
cluster-space toolkit rather than a single-use script.

%==============================================================================
% REFERENCES
% Plain integer \bibitem keys, IEEEtran style, no BibTeX.
%   [verified] = confirmed citation.   [FILL] = complete from your own records.
%==============================================================================
\begin{thebibliography}{99}

\bibitem{1} % [verify vol/no/pages against the paper]
C. A. Kitts and I. Mas, ``Cluster space specification and control of mobile
multirobot systems,'' \emph{IEEE/ASME Trans. Mechatronics}, vol.~14, no.~2,
pp.~207--218, 2009.

\bibitem{2} % [FILL] cluster-space follow-on (dynamic / partitioned control)
I. Mas and C. A. Kitts, ``\fillnote{title},'' \fillnote{venue, vol., pp., year}.

\bibitem{3} % [FILL] the 2016 SCU toolbox/thesis that the depth-quadratic
            % assembly and the flagged rotation-product recomputation come from
\fillnote{Author}, ``\fillnote{title},'' \fillnote{M.S./Ph.D.} thesis, Santa
Clara University, 2016.

\bibitem{4} % [verified]
P. Corke and J. Haviland, ``Not your grandmother's toolbox: the Robotics Toolbox
reinvented for Python,'' in \emph{Proc. IEEE Int. Conf. Robotics and Automation
(ICRA)}, 2021, pp.~11357--11363.

\bibitem{5} % [verified]
P. I. Corke, ``A robotics toolbox for MATLAB,'' \emph{IEEE Robot. Autom. Mag.},
vol.~3, no.~1, pp.~24--32, 1996.

\bibitem{6} % [verified] Procrustes / least-squares orientation alignment
S. Umeyama, ``Least-squares estimation of transformation parameters between two
point patterns,'' \emph{IEEE Trans. Pattern Anal. Mach. Intell.}, vol.~13, no.~4,
pp.~376--380, 1991.

\bibitem{7} % [verified] Sherman-Morrison / matrix computations
G. H. Golub and C. F. Van Loan, \emph{Matrix Computations}, 4th ed. Baltimore,
MD, USA: Johns Hopkins Univ. Press, 2013.

\bibitem{8} % [FILL] a formation-control / virtual-structure anchor for the intro
\fillnote{Author(s)}, ``\fillnote{virtual structure or formation control ref},''
\fillnote{venue, year}.

\bibitem{9} % [FILL] your own conference paper, for self-citation of the testbed
C. Waight \emph{et al.}, ``A functional test bed \fillnote{full title},''
\fillnote{venue, year}.

\end{thebibliography}

\end{document}
